\section{Elliptic function conventions}
\label{app:elliptic}
\subsection{Theta and elliptic gamma functions}
\label{app:theta-gamma}
We use theta functions with quasi-periods $\pi,\pi\tau$, nome
$q=e^{\ii\pi\tau}$, and $Q=q^2$. Their
product definitions are
\begin{equation}
  \begin{aligned}
    \vartheta_1(a\mid\tau)
    &=2q^{1/4}\sin a\prod_{m=1}^{\infty}
    (1-q^{2m})(1-2q^{2m}\cos2a+q^{4m}),\\
    \vartheta_2(a\mid\tau)
    &=2q^{1/4}\cos a\prod_{m=1}^{\infty}
    (1-q^{2m})(1+2q^{2m}\cos2a+q^{4m}),\\
    \vartheta_3(a\mid\tau)
    &=\prod_{m=1}^{\infty}
    (1-q^{2m})(1+2q^{2m-1}\cos2a+q^{4m-2}),\\
    \vartheta_4(a\mid\tau)
    &=\prod_{m=1}^{\infty}
    (1-q^{2m})(1-2q^{2m-1}\cos2a+q^{4m-2}),\\
    \eta(\tau)&=q^{1/12}\prod_{m=1}^{\infty}(1-q^{2m}).
  \end{aligned}
  \label{eq:theta-products}
\end{equation}
The normalisation aligns with \cite{WhittakerWatson,Lawden,DLMF}. The odd theta
function obeys
\begin{equation}
  \vartheta_1(a+\pi)=-\vartheta_1(a),\qquad
  \vartheta_1(a+\pi\tau)=-q^{-1}e^{-2\ii a}\vartheta_1(a).
  \label{eq:theta-periods}
\end{equation}
Under half-period shifts we have
\begin{equation}
  \begin{gathered}
    \vartheta_1(a+\pi/2)=\vartheta_2(a),\qquad
    \vartheta_4(a+\pi/2)=\vartheta_3(a),\\
    \vartheta_1(a\pm\pi\tau/2)
    =\pm\ii q^{-1/4}e^{\mp\ii a}\vartheta_4(a),\qquad
    \vartheta_4(a\pm\pi\tau/2)
    =\pm\ii q^{-1/4}e^{\mp\ii a}\vartheta_1(a).
  \end{gathered}
  \label{eq:half-periods}
\end{equation}
Finally we have the duplication identity
\begin{equation}
  \vartheta_1'(0)\vartheta_1(2a)
  =2\vartheta_1(a)\vartheta_2(a)\vartheta_3(a)\vartheta_4(a),
  \qquad
  \vartheta_1'(0)=\vartheta_2(0)\vartheta_3(0)\vartheta_4(0).
  \label{eq:theta-duplication}
\end{equation}

We also introduce the elliptic $Q$-theta functions
\begin{equation}
  \Theta_Q(z)=(z;Q)_\infty(Q/z;Q)_\infty,\qquad
  A_Q=(Q;Q)_\infty.
  \label{eq:multiplicative-theta}
\end{equation}
With $z=e^{2\ii a}$ these are related to the previous set of theta functions as
follows
\begin{equation}
  \begin{gathered}
    \vartheta_1(a)=\ii q^{1/4}A_Qe^{-\ii a}\Theta_Q(z),
    \qquad \vartheta_4(a)=A_Q\Theta_Q(qz).
  \end{gathered}
  \label{eq:jacobi-dictionary}
\end{equation}

For two nomes with $0<|p|,|Q|<1$, the elliptic gamma function is
\begin{equation}
  \Gamma_{p,Q}(z)
  =\prod_{j,k\geq0}
  \frac{1-p^{j+1}Q^{k+1}/z}{1-p^jQ^kz},\qquad z\in\mathbb C^*.
  \label{eq:elliptic-gamma}
\end{equation}
It is meromorphic, with possible poles at $p^{-j}Q^{-k}$ and
zeros at $p^{j+1}Q^{k+1}$, $j,k\geq0$.
We use the standard notation where a $\pm$ superscript refers to products of the
positive and negative power. For example
$\Gamma_{p,Q}(v^{\pm1}z^{\pm1})
  =\Gamma_{p,Q}(vz)\Gamma_{p,Q}(v/z)
  \Gamma_{p,Q}(z/v)\Gamma_{p,Q}(1/(vz))$.
Useful identities are
\begin{equation}
  \begin{gathered}
    \Gamma_{p,Q}(z)\Gamma_{p,Q}(pQ/z)=1,\qquad
    \Gamma_{p,Q}(pz)=\Theta_Q(z)\Gamma_{p,Q}(z),\qquad
    \Gamma_{p,Q}(Qz)=\Theta_p(z)\Gamma_{p,Q}(z).
  \end{gathered}
  \label{eq:gamma-identities}
\end{equation}

\subsection{Jacobi elliptic functions and Fourier series}
The theta/Jacobi correspondence is
\begin{equation}
  \sn s=\frac{\vartheta_3(0)}{\vartheta_2(0)}
  \frac{\vartheta_1(a)}{\vartheta_4(a)},\quad
  \cn s=\frac{\vartheta_4(0)}{\vartheta_2(0)}
  \frac{\vartheta_2(a)}{\vartheta_4(a)},\quad
  \dn s=\frac{\vartheta_4(0)}{\vartheta_3(0)}
  \frac{\vartheta_3(a)}{\vartheta_4(a)},\qquad s=\frac{2K}{\pi}a.
\end{equation}
$K(k)$ is the complete elliptic integral of the first kind and
$K'(k)=K(\sqrt{1-k^2})$ is the complementary integral. Beyond the standard
relations $\sn^2s+\cn^2s=1$ and $k^2\sn^2s+\dn^2s=1$, we need
\begin{equation}
  \begin{aligned}
    \sn(2s)&=\frac{2\sn s\cn s\dn s}{1-k^2\sn^4s},\\
    \sn(s+t)\sn(s-t)&=
    \frac{\sn^2s-\sn^2t}{1-k^2\sn^2s\sn^2t},\\
    \cn(s-t)-\cn(s+t)&=
    \frac{2\sn s\sn t\dn s\dn t}{1-k^2\sn^2s\sn^2t}.
  \end{aligned}
\end{equation}
For $L=\pi K'/K$, the relevant Fourier series are
\begin{align}
  \sn s&=\frac\pi{kK}\sum_{r\geq0}
  \csch\!\left(\frac{(2r+1)L}{2}\right)
  \sin\!\left(\frac{(2r+1)\pi s}{2K}\right),\label{eq:sn-fourier}\\
  \cn s&=\frac\pi{kK}\sum_{r\geq0}
  \sech\!\left(\frac{(2r+1)L}{2}\right)
  \cos\!\left(\frac{(2r+1)\pi s}{2K}\right),\label{eq:cn-fourier}\\
  \dn s&=\frac\pi{2K}+\frac\pi K\sum_{r\geq1}
  \sech(rL)\cos\!\left(\frac{r\pi s}{K}\right).
  \label{eq:dn-fourier}
\end{align}
